% 《大学物理学学习指导》第5版，自测题11“量子物理”
% 选择11.1.1—11.1.12，印刷230—231，PDF物理238—239页。
\subsection{自测题11：光电、康普顿、物质波与量子态选择题}

\sourcetest{11.1.1}{230}{238}
\problemtext 金属光电逸出需要能量$eU_0$，
单色入射光真空波长$\lambda$必须满足何种阈值条件？
\workings
发生外光电效应的条件是
$hc/\lambda\ge eU_0$，于是
\[
\boxed{\lambda\le\frac{hc}{eU_0}}.
\]
选\textbf{A}。
\crosscheck 波长越短光子能量越大；
不应把阈值不等式方向反写。

\sourcetest{11.1.2}{230}{238}
\problemtext 同一金属、同一实验中，
照射波长由400nm减至300nm，
问遏止电压如何变化。
\workings
爱因斯坦光电方程$eU_a=hc/\lambda-W$，
同种金属作差消去逸出功，
\[
\Delta U_a=\frac{hc}{e}
\left(\frac1{300\,\mathrm{nm}}-\frac1{400\,\mathrm{nm}}\right)
\simeq\boxed{+1.033\,\mathrm V}.
\]
故遏止电压\textbf{增大约1.035V}，选\textbf{D}。
\crosscheck 题给选项按约两位小数近似；
不是随波长缩短而降低电压。

\sourcetest{11.1.3}{230}{238}
\problemtext 保持真空光电管阳极—阴极电势差不变，
仅增大照射同一阴极的同频单色光强，
判断逸出电子的最大初动能及到阳极时最大动能变化。
\workings
光强增大增加单位时间入射光子数与可能光电子数，
但每个光子能量$h\nu$保持不变，
$\boxed{K_{0,\max}=h\nu-W}$不变。
静电势差不变，电势做功$e|\Delta V|$也不变，
到阳极的最大动能同样不变。
选\textbf{D（均不变）}。
\crosscheck 电流大小可能增大，
不能将光电子数增加误当作单个光电子能量增加。

\sourcetest{11.1.4}{230}{238}
\problemtext X射线散射实验中的散射光可能包含哪些波长成分，
康普顿波长改变是否依赖材料？
\workings
可出现与入射波长相同的弹性（相干）散射分量，
以及波长增长的非弹性康普顿散射分量。
理想自由电子的康普顿位移
\[
\boxed{\Delta\lambda=\frac{h}{m_ec}(1-\cos\theta)}
\]
只依散射角$\theta$和电子静质量，
不因材料的化学名称而改变该基本位移关系。
选\textbf{B}。
\crosscheck 各散射分量的\emph{强度}会明显依赖材料、
束缚结构和入射能量；
不能把“波长位移的理想公式与材料无关”
误写成全部散射现象与材料无关。

\sourcetest{11.1.5}{230}{238}
\problemtext 康普顿散射中散射后光子波长是原来的1.2倍，
求散射后光子能量与反冲电子动能之比。
\workings
光子能量反比波长：
$E'_\gamma=E_\gamma/1.2=5E_\gamma/6$，
自由静止电子反冲动能
$K=E_\gamma-E'_\gamma=E_\gamma/6$，故
\[
\boxed{E'_\gamma/K=5}.
\]
选\textbf{D}。
\crosscheck 题问的是\emph{散射后}光子能量，
不是初始能量与反冲动能之比6。

\sourcetest{11.1.6}{230--231}{238--239}
\problemtext 康普顿散射中原先静止电子被反冲到$0.6c$，
求其动能是静止能量的几倍。
\workings
$\gamma=(1-0.6^2)^{-1/2}=1.25$，
\[
\boxed{\frac{K}{m_ec^2}=\gamma-1=0.25}.
\]
选\textbf{D}。
\crosscheck $0.5m_ev^2=0.18m_ec^2$是经典动能，
在本题速度下误差不可忽略。

\sourcetest{11.1.7}{231}{239}
\problemtext 玻尔氢原子模型中，电子第一与第三轨道速度比$v_1/v_3$。
\workings
玻尔氢原子第$n$轨道角动量量子化
$m_ev_nr_n=n\hbar$，向心库仑力与轨道关系联立
得到$v_n\propto1/n$，故
\[
\boxed{v_1/v_3=3}.
\]
选\textbf{C}。
\crosscheck 此为玻尔圆轨道模型，不能作为现代量子力学
电子真实经典圆轨迹的陈述。

\sourcetest{11.1.8}{231}{239}
\problemtext 氢原子从基态激发到某态需10.19eV；
若又从能量$-0.85$eV的状态跃迁到该激发态，
发射光子能量是多少？
\workings
基态$E_1=-13.6$eV，
加10.19eV得到目标态$E_f=-3.41$eV
（按13.6eV标准近似为$n=2$，$E_2=-3.4$eV）。
从$-0.85$eV（$n=4$）到$n=2$，
发射能量
\[
\boxed{E_\gamma=(-0.85)-(-3.40)=2.55\,\mathrm{eV}
\simeq2.56\,\mathrm{eV}}.
\]
选\textbf{A}。
\crosscheck 原题10.19eV只保留四位数字，
若始终使用它得到2.56eV；
采用精确玻尔13.6eV能级表则是2.55eV，
必须注明数据精度来源。

\sourcetest{11.1.9}{231}{239}
\problemtext 两种不同质量的粒子德布罗意波长相等，
哪一项物理量必然相等？
\workings
德布罗意关系
$\lambda=h/p$给
\[
\boxed{p_1=p_2=h/\lambda},
\]
故是\textbf{动量大小相同}，选\textbf{A}。
\crosscheck 速度、动能和总能量不必相同；
这里无须假定二者具有相同静质量。

\sourcetest{11.1.10}{231}{239}
\problemtext 原图11.1.1显示四幅沿$x$方向的粒子波包(A)—(D)，
问哪种波包有最精确的动量。
\workings
位置分布空间越宽，波数分布越窄，
则$\Delta p_x=\hbar\,\Delta k$越小。
在原图四种相似的局域波包中，
(A)具有最长、最多周期的波列，
波数相对散布最窄，故选\textbf{A}。
\crosscheck \textbf{图形比较为定性：}
原图并未给各波包归一化强度或定量谱宽；
“周期数最多”只是判断较长相干波列动量分布
可能较窄的启发性证据。若需严密排序，
应对四幅波形给出明确解析表达并计算傅里叶谱标准差。

\sourcetest{11.1.11}{231}{239}
\problemtext 波长500nm光子的波长不确定度为$10^{-4}$nm；
按\emph{原题指定的不确定关系}
$\Delta p_x\,\Delta x\ge h$求$x$位置至少多不确定。
\workings
$p=h/\lambda$，小不确定度下
$\Delta p\simeq h\Delta\lambda/\lambda^2$。
于是按原题给定的$h$形式，
\[
\boxed{\Delta x\gtrsim\frac h{\Delta p}
=\frac{\lambda^2}{\Delta\lambda}
=\frac{(500\,\mathrm{nm})^2}{10^{-4}\,\mathrm{nm}}
=2.5\,\mathrm m=250\,\mathrm{cm}}.
\]
选\textbf{C}。
\crosscheck \textbf{不等式常数约定：}
规范位置—动量标准差下界为$\hbar/2=h/(4\pi)$。
同时光子的严格位置算符有理论细节；
本条只是在\emph{题目所给波包近似关系}下计算，
不将其当作普适的光子严格定位定理。

\sourcetest{11.1.12}{231}{239}
\problemtext 氢原子L壳层（$n=2$）电子的四量子数
$(n,l,m_l,m_s)$，哪组满足允许值？
\workings
当$n=2$时$l=0,1$；
$l=1$时$m_l=-1,0,1$；
电子$m_s=\pm1/2$。
原选项中满足以上规则的是
\[
\boxed{(n,l,m_l,m_s)=(2,1,-1,1/2)},
\]
选\textbf{B}。
\crosscheck 不能把$n=1$的K壳层或$n=3$的M壳层组合
计入L壳层；$l=0$时$m_l$只能为0。
